\documentclass{article}

\usepackage[utf8]{inputenc}
\usepackage{amsfonts}
\usepackage{amssymb}
\usepackage{amsmath}
\usepackage{amsthm}
\usepackage{enumitem}
\usepackage{bbold}
\usepackage{bm}
\usepackage{graphicx}
\usepackage{color}
\usepackage{hyperref}
\usepackage[margin=2.5cm]{geometry}

\begin{document}

% ==============================================================================

\title{\Large{INFO8006: Project 2 - Report}}

\vspace{1cm}

\author{\small{\bf phoomchit samoban - b6735128}}

\maketitle

% ==============================================================================

\section{Bayes filter}

\begin{enumerate}[label=\alph*.,leftmargin=*]

    \item
    The rusty sensor provides a noisy measurement of the Manhattan distance
    between Pacman and a ghost.

    Let $X_t=(x,y)$ denote the ghost position at time $t$, and let
    $d(X_t,P_t)$ denote the Manhattan distance between the ghost and Pacman.
    The sensor noise is generated from a binomial random variable

    \[
    K \sim \mathrm{Binomial}(n,p),
    \]

    where

    \[
    p=0.5
    \]

    and

    \[
    n=
    \left\lfloor
    \frac{\text{sensor variance}}{p(1-p)}
    \right\rfloor.
    \]

    The observed evidence is therefore

    \[
    E_t=d(X_t,P_t)+(K-np).
    \]

    For a possible ghost position $x$, the sensor model is

    \[
    P(E_t=e\mid X_t=x)
    =
    \Pr\left(
    K=e-d(x,P_t)+np
    \right).
    \]

    The probability is zero when the required value of $K$ is outside the
    valid range from $0$ to $n$.

    \item
    The transition model describes how a ghost moves from its current
    position to a legal neighbouring position.

    Let $d$ be the current Manhattan distance between the ghost and Pacman,
    and let $d'$ be the distance after a possible move. The probability of
    a legal action is proportional to

    \[
    w(a)=
    \begin{cases}
    \lambda, & d' \geq d,\\
    1, & d' < d.
    \end{cases}
    \]

    After assigning the weights, they are normalized over all legal
    non-stop actions.

    The parameter $\lambda$ depends on the ghost type:

    \[
    \lambda=
    \begin{cases}
    1, & \text{confused},\\
    2, & \text{afraid},\\
    8, & \text{scared}.
    \end{cases}
    \]

    Therefore, a confused ghost moves uniformly, while afraid and scared
    ghosts prefer movements that maintain or increase their distance from
    Pacman.

\end{enumerate}

% ==============================================================================

\section{Implementation}

\begin{enumerate}[label=\alph*.,leftmargin=*]

    \item
    The Bayes filter is implemented by maintaining a separate belief
    distribution for each ghost.

    At each time step, the belief is updated using three steps.

    First, the prediction step applies the transition model:

    \[
    \overline{Bel}_t(x)
    =
    \sum_{x'}
    P(x\mid x')Bel_{t-1}(x').
    \]

    Second, the measurement update applies the sensor model:

    \[
    \widetilde{Bel}_t(x)
    =
    P(E_t\mid x)\overline{Bel}_t(x).
    \]

    Finally, the result is normalized:

    \[
    Bel_t(x)
    =
    \frac{
    \widetilde{Bel}_t(x)
    }{
    \sum_y \widetilde{Bel}_t(y)
    }.
    \]

    The current Pacman position is used in the sensor and transition
    calculations. If a ghost has been eaten, its belief is set to zero.
    The implementation also handles the case where the updated belief has
    total probability zero by falling back to a uniform distribution.

\end{enumerate}

% ==============================================================================

\section{Experiment}

\begin{enumerate}[label=\alph*.,leftmargin=*]

    \item
    We use entropy as a measure of uncertainty in the belief distribution:

    \[
    H(B)
    =
    -\sum_x B(x)\log_2 B(x).
    \]

    A lower entropy means that the probability is concentrated on fewer
    positions, while a higher entropy represents greater uncertainty.

    \item
    To measure the quality of the belief state, we use the expected
    Manhattan distance between the predicted ghost position and the true
    ghost position:

    \[
    Q(B)
    =
    \sum_x
    B(x)
    d_{\mathrm{Manhattan}}(x,x_{\mathrm{true}}).
    \]

    Lower values indicate that the belief distribution is more concentrated
    around the true ghost position.

    \item
    The implementation was tested on the two required layouts using one
    hidden ghost and each of the three ghost types.

    \begin{center}
    \begin{tabular}{lccc}
    \hline
    Layout & Ghost type & Score & Result \\
    \hline
    large\_filter & confused & 698 & Win \\
    large\_filter & afraid & 652 & Win \\
    large\_filter & scared & 518 & Win \\
    large\_filter\_walls & confused & 698 & Win \\
    large\_filter\_walls & afraid & 645 & Win \\
    large\_filter\_walls & scared & 659 & Win \\
    \hline
    \end{tabular}
    \end{center}

    All six tests completed successfully without the Pacman agent losing
    the game.

    The measured computation times were approximately 0.012 seconds to
    1.69 seconds for the tested cases.

    These results show that the implementation can successfully track hidden
    ghosts in both layouts for confused, afraid, and scared ghost behaviour.

    \item
    The transition parameter controls how strongly the ghost prefers actions
    that maintain or increase its distance from Pacman.

    When the parameter is larger, actions that increase or maintain the
    distance receive a larger probability compared with actions that reduce
    the distance. Therefore, increasing the transition parameter produces
    stronger fleeing behaviour.

    \item
    Sensor variance determines the amount of noise in the distance
    observation. A smaller sensor variance produces more accurate evidence,
    allowing the belief distribution to become more concentrated around
    likely ghost positions.

    A larger sensor variance produces noisier evidence and therefore causes
    greater uncertainty in the belief distribution. Consequently, entropy
    is expected to increase as the sensor variance increases, while belief
    quality is expected to decrease.

    \item
    The Pacman controller uses the current Pacman position, legal actions,
    and the belief distributions produced by the Bayes filter.

    The belief distributions of ghosts that have not been eaten are combined
    into a risk map. The most likely ghost position is then used as a target.
    The controller evaluates legal Pacman actions and avoids positions with
    sufficiently high ghost probability whenever possible.

    When a safe path is available, Pacman moves toward the estimated ghost
    position. If all possible moves are considered unsafe, the controller
    uses a fallback strategy and selects the action with the lowest estimated
    ghost risk.

    \item \textbf{Leave empty.}

\end{enumerate}

% ==============================================================================

\end{document}